\section{Conclusioni}
\label{sec:conclusioni}

\subsection{Sintesi dei risultati}

Il progetto ha sviluppato e confrontato tre configurazioni di uno sciame eterogeneo per un problema di raccolta e consegna su griglia. I risultati dimostrano con chiarezza una gerarchia di efficienza:

\begin{center}
\begin{tabular}{lcc}
\toprule
\textbf{Configurazione} & \textbf{Tick (Mappa A)} & \textbf{Tick (Mappa B)} \\
\midrule
Exploration & 251 & 153 \\
Collection  & 129 & 117 \\
With Relay  & 133 & 117 \\
\bottomrule
\end{tabular}
\end{center}

La configurazione \emph{Exploration} si rivela la meno efficiente in assoluto: il surplus di agenti esplorativi non compensa la carenza di capacità operativa, e la comunicazione scout-collector basata unicamente sull'incontro casuale introduce latenze elevate. Il divario è particolarmente evidente su Mappa~A (251 vs.\ 129 tick), dove la mappa aperta avrebbe potuto favorire incontri più frequenti ma non riesce a colmare il gap nel parallelismo della raccolta.

\emph{Collection} e \emph{With Relay} convergono a performance simili, ma con meccanismi diversi: la prima punta sul parallelismo grezzo dei Collector, la seconda sull'intelligenza del routing informativo.

\subsection{Il Relay come soluzione architetturale}

Il contributo principale della configurazione \emph{With Relay} non è la superiorità numerica bruta, ma l'\textbf{architettura di comunicazione più robusta}. Il Relay trasforma la topologia della rete da un grafo casuale e dipendente dalle traiettorie degli agenti a una struttura con un nodo centrale garantito.

Questa robustezza architetturale emerge con forza su Mappa~B: nonostante abbia un Collector in meno rispetto a \emph{Collection} (2 vs.\ 3), \emph{With Relay} eguaglia le performance perché l'informazione sugli oggetti viene propagata più velocemente e con meno dipendenza dalla fortuna delle traiettorie. Il Relay ``garantisce'' che l'informazione fluisca, indipendentemente da dove si trovino Scout e Collector in quel momento.

\subsection{Dipendenza dalla struttura della mappa}

Il risultato forse più importante del progetto è la dimostrazione che \textbf{non esiste una configurazione universalmente ottimale}: la scelta dipende dalla topologia della mappa.

\begin{itemize}
  \item \textbf{Mappa aperta (pochi ostacoli)}: la comunicazione diretta è già efficiente perché gli agenti si incontrano spesso. Il Relay aggiunge overhead organizzativo senza vantaggio proporzionale. Conviene investire nel parallelismo della raccolta (\emph{Collection}).
  
  \item \textbf{Mappa complessa (molti ostacoli, corridoi stretti)}: la comunicazione diretta è inaffidabile perché i percorsi obbligati separano Scout e Collector. Il Relay -- posizionato nel centro, passaggio naturale per tutti i corridoi -- diventa indispensabile per la propagazione rapida delle informazioni.
  
  \item \textbf{Mappa estremamente complessa} (ipotetica): ci si aspetta che il vantaggio del Relay cresca ulteriormente, al punto da giustificare anche due Relay o una rete di relay distribuiti.
\end{itemize}

Questa dipendenza dalla struttura è un risultato fondamentale nella letteratura di swarm robotics: le strategie di coordinamento non sono universali ma devono essere \emph{co-progettate con l'ambiente di deployment}. Il Relay è, in questo senso, una risposta adattiva alla sfida comunicativa imposta dalla topologia.

\subsection{Possibili sviluppi futuri}

\begin{itemize}
  \item \textbf{Relay adattivo}: invece di una traiettoria fissa a diamante, il Relay potrebbe adottare un percorso calcolato dinamicamente in base alla distribuzione corrente degli agenti e degli oggetti noti.
  \item \textbf{Meccanismi stigmergici}: deposito di feromoni virtuali nelle celle per guidare l'esplorazione verso zone non ancora visitate, riducendo l'overhead del calcolo esplicito delle frontiere.
  \item \textbf{Topologie multi-relay}: su mappe di dimensione maggiore o con frammentazione più elevata, una rete di più Relay distribuiti potrebbe garantire copertura comunicativa completa.
  \item \textbf{Batteria eterogenea}: assegnare batterie differenti ai vari ruoli (maggiore per il Relay che deve pattugliare a lungo, minore per i Collector che restano vicini ai magazzini) per ottimizzare il budget energetico complessivo.
\end{itemize}
